\begin{frame}{``왜 2012년이었을까'' --- 두 번째 축: 데이터}
  \begin{itemize}
    \item ImageNet: \textbf{약 140만 장}의 이미지에 \textbf{1000개 클래스}
      라벨을 붙인 데이터셋 (2009년 공개).
    \item 중요한 점: 라벨 대부분이 \textbf{일반인이 직접 매긴 것} ---
      인터넷과 크라우드소싱(아마존 메카니컬 터크) 플랫폼이 없었다면,
      100만 장 넘는 사진마다 ``고양이''·``셔터''를 붙이는 비용이
      \textbf{연구비를 압도}했을 것이다.
  \end{itemize}
  \vskip0.6em
  \begin{block}{``140만 장''과 ``2장 GPU''의 동시 출현}
    AlexNet의 승리는 두 축이 \textbf{동시에} 갖춰진 2012년에만 가능했던 것 ---
    어느 하나만 있었어도 그 해의 결과물은 달라졌다.
  \end{block}
  \vskip0.4em
  이번 학기 내내 등장하는 모든 모델이 이 ``두 축'' 위에 선다.
\end{frame}
